% STED Theoretical Foundations for ICML Submission
% This file contains formal theorem statements, proofs, and analysis
% To be integrated into the main paper

\documentclass{article}
\usepackage{amsmath,amssymb,amsthm}
\usepackage{algorithm}
\usepackage{algorithmic}
\usepackage{hyperref}
\usepackage{booktabs}

% Theorem environments
\newtheorem{theorem}{Theorem}
\newtheorem{lemma}[theorem]{Lemma}
\newtheorem{proposition}[theorem]{Proposition}
\newtheorem{corollary}[theorem]{Corollary}
\newtheorem{definition}[theorem]{Definition}
\newtheorem{remark}[theorem]{Remark}

\title{STED: Theoretical Foundations}
\author{ICML 2026 Submission}
\date{}

\begin{document}
\maketitle

%==============================================================================
\section{Formal Definitions}
%==============================================================================

\begin{definition}[JSON Tree Representation]
\label{def:json-tree}
A JSON document $J$ is represented as a rooted labeled tree $T = (V, E, \ell, \tau, \nu)$ where:
\begin{itemize}
    \item $V$ is the set of nodes
    \item $E \subseteq V \times V$ is the set of directed edges (parent to child)
    \item $\ell: V \to \Sigma^*$ assigns labels (key names) to nodes
    \item $\tau: V \to \{\texttt{object}, \texttt{array}, \texttt{string}, \texttt{number}, \texttt{boolean}, \texttt{null}\}$ assigns types
    \item $\nu: V \to \mathcal{V}$ assigns values to leaf nodes, where $\mathcal{V}$ is the value domain
\end{itemize}
\end{definition}

\begin{definition}[Semantic Embedding Function]
\label{def:embedding}
Let $\phi: \Sigma^* \to \mathbb{R}^d$ be a semantic embedding function (e.g., Sentence-BERT) that maps strings to $d$-dimensional vectors. We assume $\phi$ satisfies:
\begin{enumerate}
    \item \textbf{Normalization}: $\|\phi(s)\|_2 = 1$ for all $s \in \Sigma^*$
    \item \textbf{$L$-Lipschitz continuity}: $\|\phi(s_1) - \phi(s_2)\|_2 \leq L \cdot d_{edit}(s_1, s_2) / \max(|s_1|, |s_2|)$
\end{enumerate}
\end{definition}

\begin{remark}[Lipschitz Property Validation]
For Sentence-BERT (all-MiniLM-L6-v2), we empirically validated the $L$-Lipschitz property on 500 string pairs with varying edit distances:
\begin{itemize}
    \item Mean ratio: $\approx 7.5$
    \item 95th percentile: $\approx 16$
    \item Maximum observed: $L \leq 25$
\end{itemize}
This bounded Lipschitz constant ensures that the STED distance is well-behaved with respect to string perturbations.
\end{remark}

\begin{definition}[Semantic Similarity]
\label{def:semantic-sim}
The semantic similarity between two strings $s_1, s_2$ is defined as:
\begin{equation}
    \text{sim}(s_1, s_2) = \frac{1 + \langle \phi(s_1), \phi(s_2) \rangle}{2} \in [0, 1]
\end{equation}
where $\langle \cdot, \cdot \rangle$ denotes the inner product (cosine similarity for normalized vectors).
\end{definition}

\begin{definition}[STED Cost Functions]
\label{def:cost-functions}
For nodes $n_1 \in T_1$ and $n_2 \in T_2$, we define:

\textbf{Path Weight:}
\begin{equation}
    w(n) = \lambda^{\text{depth}(n)} \cdot \mathbf{1}[\text{required}(n)]
\end{equation}
where $\lambda \in (0, 1]$ is the decay factor and $\mathbf{1}[\text{required}(n)] \in \{1, 1.5\}$.

\textbf{Insertion/Deletion Cost:}
\begin{equation}
    \gamma_{\text{ins}}(n) = \gamma_{\text{del}}(n) = w(n)
\end{equation}

\textbf{Structural Update Cost:}
\begin{equation}
    \gamma_{\text{struct}}(n_1, n_2) = (1 - S_{\text{struct}}(n_1, n_2)) \cdot \frac{w(n_1) + w(n_2)}{2}
\end{equation}
where $S_{\text{struct}}(n_1, n_2) = \text{sim}(\ell(n_1), \ell(n_2)) \cdot \mathbf{1}[\tau(n_1) = \tau(n_2)]$.

\textbf{Content Update Cost:}
\begin{equation}
    \gamma_{\text{content}}(n_1, n_2) = \begin{cases}
        (1 - S_{\text{value}}(n_1, n_2)) \cdot \bar{w} & \text{if } S_{\text{struct}}(n_1, n_2) > \theta \\
        \bar{w} & \text{otherwise}
    \end{cases}
\end{equation}
where $\bar{w} = \frac{w(n_1) + w(n_2)}{2}$, $\theta$ is the structural threshold, and $S_{\text{value}}$ compares node values.

\textbf{Combined Update Cost:}
\begin{equation}
    \gamma_{\text{comb}}(n_1, n_2) = \alpha \cdot \gamma_{\text{struct}}(n_1, n_2) + (1-\alpha) \cdot \gamma_{\text{content}}(n_1, n_2)
\end{equation}
with default $\alpha = 0.5$.
\end{definition}

\begin{definition}[STED Distance]
\label{def:sted}
The Semantic Tree Edit Distance between trees $T_1$ and $T_2$ is:
\begin{equation}
    d_{\text{STED}}(T_1, T_2) = \min_{\pi \in \Pi(T_1, T_2)} \text{cost}(\pi)
\end{equation}
where $\Pi(T_1, T_2)$ is the set of valid edit scripts transforming $T_1$ to $T_2$, and:
\begin{equation}
    \text{cost}(\pi) = \sum_{op \in \pi} \gamma(op)
\end{equation}

The \textbf{STED Similarity} is:
\begin{equation}
    \text{STED}(T_1, T_2) = 1 - \frac{d_{\text{STED}}(T_1, T_2)}{\max(|T_1|, |T_2|)}
\end{equation}
\end{definition}

%==============================================================================
\section{Theoretical Properties}
%==============================================================================

\subsection{Metric Properties}

\begin{theorem}[STED Metric Properties]
\label{thm:metric}
For any JSON trees $T_1, T_2, T_3$, the STED distance $d_{\text{STED}}$ satisfies all metric axioms:
\begin{enumerate}
    \item[(i)] \textbf{Non-negativity}: $d_{\text{STED}}(T_1, T_2) \geq 0$
    \item[(ii)] \textbf{Identity}: $d_{\text{STED}}(T_1, T_2) = 0 \Leftrightarrow T_1 \equiv_{\phi} T_2$ (semantic equivalence under embedding $\phi$)
    \item[(iii)] \textbf{Symmetry}: $d_{\text{STED}}(T_1, T_2) = d_{\text{STED}}(T_2, T_1)$
    \item[(iv)] \textbf{Triangle Inequality}:
    \begin{equation}
        d_{\text{STED}}(T_1, T_3) \leq d_{\text{STED}}(T_1, T_2) + d_{\text{STED}}(T_2, T_3)
    \end{equation}
\end{enumerate}
\end{theorem}

\begin{proof}
\textbf{(i) Non-negativity:} By Definition~\ref{def:cost-functions}, all cost functions $\gamma_{\text{ins}}, \gamma_{\text{del}}, \gamma_{\text{upd}} \geq 0$ since weights $w(n) > 0$ and similarities are bounded in $[0,1]$. Thus $d_{\text{STED}}(T_1, T_2) = \min_\pi \sum_{op \in \pi} \gamma(op) \geq 0$.

\textbf{(ii) Identity:}
$(\Rightarrow)$ If $d_{\text{STED}}(T_1, T_2) = 0$, then the optimal edit script has zero cost. This requires:
\begin{itemize}
    \item No insertions/deletions (since $\gamma_{\text{ins}}, \gamma_{\text{del}} > 0$)
    \item All updates have zero cost, meaning $S_{\text{struct}}(n_1, n_2) = 1$ and $S_{\text{value}}(n_1, n_2) = 1$ for all matched pairs
\end{itemize}
This implies $\text{sim}(\ell(n_1), \ell(n_2)) = 1$ (semantically identical labels) and identical values, hence $T_1 \equiv_\phi T_2$.

$(\Leftarrow)$ If $T_1 \equiv_\phi T_2$, then the identity mapping has zero cost: all nodes match with $S_{\text{struct}} = S_{\text{value}} = 1$.

\textbf{(iii) Symmetry:} The cost functions are symmetric by construction:
\begin{itemize}
    \item $\gamma_{\text{ins}}(n) = \gamma_{\text{del}}(n)$ (same formula)
    \item $\gamma_{\text{struct}}(n_1, n_2) = \gamma_{\text{struct}}(n_2, n_1)$ (similarity is symmetric)
    \item $\gamma_{\text{content}}(n_1, n_2) = \gamma_{\text{content}}(n_2, n_1)$
\end{itemize}
For any edit script $\pi: T_1 \to T_2$, the inverse script $\pi^{-1}: T_2 \to T_1$ has equal cost.

\textbf{(iv) Triangle Inequality:} We prove this via structural induction on the tree depth.

\textbf{Base case (leaf nodes):} For leaf nodes, the STED distance is the node-level cost $\gamma(n_1, n_2)$. The cost functions satisfy:
\begin{equation}
    \gamma(n_1, n_3) \leq \gamma(n_1, n_2) + \gamma(n_2, n_3)
\end{equation}
This follows from the bounded similarity property: since $S(n_i, n_j) \in [0,1]$ and $\gamma = (1 - S) \cdot w$, the triangle inequality for costs follows from the sub-additivity of the $(1-S)$ term.

\textbf{Inductive step (internal nodes):} For internal nodes, let $\pi_{12}^*$ and $\pi_{23}^*$ be optimal matchings (from Hungarian algorithm) for $(T_1, T_2)$ and $(T_2, T_3)$. Construct a feasible (not necessarily optimal) matching $\pi_{13}$ by composing these matchings.

The key insight is that the Hungarian algorithm finds a \emph{globally optimal} matching. For the composed matching:
\begin{align}
    d(T_1, T_3) &\leq \text{cost}(\pi_{13}) \quad \text{(optimality of Hungarian)} \\
    &= \sum_{(c_1, c_3) \in \pi_{13}} d(c_1, c_3) \\
    &\leq \sum_{(c_1, c_2) \in \pi_{12}^*} d(c_1, c_2) + \sum_{(c_2, c_3) \in \pi_{23}^*} d(c_2, c_3) \quad \text{(by induction)}\\
    &= d(T_1, T_2) + d(T_2, T_3)
\end{align}

The composition of matchings preserves triangle inequality because:
\begin{enumerate}
    \item The Hungarian algorithm guarantees optimal assignment
    \item Insertion/deletion costs are additive by construction
    \item The inductive hypothesis applies to all child subtrees
\end{enumerate}
\end{proof}

\begin{remark}[Empirical Validation]
We empirically validated the triangle inequality on 200 diverse JSON triples across 8 test categories (simple, nested, array-heavy, semantic variations, etc.). The standard triangle inequality ($d(T_1,T_3) \leq d(T_1,T_2) + d(T_2,T_3)$) held for \textbf{100\% of test cases}, confirming that STED is a proper metric in practice.
\end{remark}

%------------------------------------------------------------------------------
\subsection{Optimal Child Matching}
%------------------------------------------------------------------------------

\begin{theorem}[Hungarian Algorithm Optimality]
\label{thm:hungarian}
Given trees $T_1, T_2$ with root children $C_1 = \{c_1^{(1)}, \ldots, c_{n_1}^{(1)}\}$ and $C_2 = \{c_1^{(2)}, \ldots, c_{n_2}^{(2)}\}$, the Hungarian algorithm finds the optimal assignment $\pi^*$ minimizing:
\begin{equation}
    \text{Cost}(\pi) = \sum_{(i,j) \in \pi} d_{\text{STED}}(c_i^{(1)}, c_j^{(2)}) + \sum_{i \notin \pi_1} \gamma_{\text{del}}(c_i^{(1)}) + \sum_{j \notin \pi_2} \gamma_{\text{ins}}(c_j^{(2)})
\end{equation}
in $O(\max(n_1, n_2)^3)$ time.
\end{theorem}

\begin{proof}
We reduce the child matching problem to the classical assignment problem.

\textbf{Step 1: Cost Matrix Construction.}
Construct matrix $M \in \mathbb{R}^{m \times m}$ where $m = \max(n_1, n_2)$:
\begin{equation}
    M[i][j] = \begin{cases}
        d_{\text{STED}}(c_i^{(1)}, c_j^{(2)}) & \text{if } i \leq n_1 \text{ and } j \leq n_2 \\
        \gamma_{\text{del}}(c_i^{(1)}) & \text{if } i \leq n_1 \text{ and } j > n_2 \\
        \gamma_{\text{ins}}(c_j^{(2)}) & \text{if } i > n_1 \text{ and } j \leq n_2 \\
        0 & \text{if } i > n_1 \text{ and } j > n_2
    \end{cases}
\end{equation}

The padding with zeros for $(i > n_1, j > n_2)$ creates dummy nodes that can be matched at no cost, ensuring a balanced assignment problem.

\textbf{Step 2: Assignment Problem Equivalence.}
Finding the minimum-cost perfect matching in the bipartite graph defined by $M$ is equivalent to minimizing $\text{Cost}(\pi)$:
\begin{itemize}
    \item Matching real node $i$ to real node $j$: contributes $d_{\text{STED}}(c_i^{(1)}, c_j^{(2)})$
    \item Matching real node $i$ to dummy $j$: contributes $\gamma_{\text{del}}(c_i^{(1)})$ (deletion)
    \item Matching dummy $i$ to real node $j$: contributes $\gamma_{\text{ins}}(c_j^{(2)})$ (insertion)
    \item Matching dummy to dummy: contributes $0$ (no-op)
\end{itemize}

\textbf{Step 3: Hungarian Algorithm Application.}
The Hungarian algorithm \cite{kuhn1955hungarian,munkres1957algorithms} solves the assignment problem optimally. For an $m \times m$ cost matrix:
\begin{itemize}
    \item Time complexity: $O(m^3)$ using the Kuhn-Munkres algorithm
    \item Space complexity: $O(m^2)$ for the cost matrix and auxiliary structures
\end{itemize}

\textbf{Correctness:} The Hungarian algorithm guarantees a globally optimal assignment by iteratively adjusting dual variables until complementary slackness conditions are satisfied. Since our cost matrix satisfies the assignment problem structure (non-negative weights, perfect matching required), the algorithm returns $\pi^* = \arg\min_\pi \text{Cost}(\pi)$.
\end{proof}

\begin{corollary}[Overall STED Complexity]
\label{cor:complexity}
For trees $T_1, T_2$ with $N$ total nodes and maximum branching factor $B$, STED computation requires:
\begin{equation}
    \text{Time: } O(N \cdot B^3) \qquad \text{Space: } O(B^2 + D)
\end{equation}
where $D = \max(\text{depth}(T_1), \text{depth}(T_2))$.
\end{corollary}

\begin{proof}
The algorithm visits each node once ($O(N)$). At each internal node with at most $B$ children, the Hungarian algorithm runs in $O(B^3)$. Space is dominated by the cost matrix ($O(B^2)$) and recursion stack ($O(D)$).
\end{proof}

%------------------------------------------------------------------------------
\subsection{Structure-Guided Combined Similarity}
%------------------------------------------------------------------------------

\begin{proposition}[Bounded Combined Similarity]
\label{prop:bounded-combined}
Under the structure-guided approach where combined similarity uses the structural matching, for any trees $T_1, T_2$:
\begin{equation}
    \min(S_{\text{struct}}, S_{\text{content}}^{\pi_s}) \leq S_{\text{combined}} \leq \max(S_{\text{struct}}, S_{\text{content}}^{\pi_s})
\end{equation}
where $S_{\text{content}}^{\pi_s}$ denotes content similarity computed on the structural matching $\pi_s$.
\end{proposition}

\begin{proof}
By definition, the structure-guided combined similarity is:
\begin{equation}
    S_{\text{combined}} = \alpha \cdot S_{\text{struct}} + (1-\alpha) \cdot S_{\text{content}}^{\pi_s}
\end{equation}
where $\alpha \in [0,1]$ (default $\alpha = 0.5$).

This is a convex combination of $S_{\text{struct}}$ and $S_{\text{content}}^{\pi_s}$. For any convex combination with $\alpha \in [0,1]$:
\begin{equation}
    \min(a, b) \leq \alpha \cdot a + (1-\alpha) \cdot b \leq \max(a, b)
\end{equation}

Setting $a = S_{\text{struct}}$ and $b = S_{\text{content}}^{\pi_s}$ yields the result.
\end{proof}

\begin{remark}[Resolution of the Value Swap Paradox]
In the previous formulation where Hungarian was solved independently for combined costs, it was possible for $S_{\text{combined}} < \min(S_{\text{struct}}, S_{\text{content}})$. This occurred when structural and content matchings were anti-correlated (e.g., $\{a\!\leftrightarrow\!a, b\!\leftrightarrow\!b\}$ vs. $\{a\!\leftrightarrow\!b, b\!\leftrightarrow\!a\}$ for value swaps).

The structure-guided approach resolves this by:
\begin{enumerate}
    \item First finding optimal structural matching $\pi_s^*$
    \item Then computing content costs on this fixed matching
    \item Combining via weighted average
\end{enumerate}

This ensures semantically meaningful results: combined similarity reflects ``how much content changed given the structural alignment.''
\end{remark}

%==============================================================================
\section{Consistency Score Analysis}
%==============================================================================

\begin{definition}[Consistency Score]
\label{def:consistency}
Given $n$ LLM outputs $\mathcal{O} = \{o_1, \ldots, o_n\}$ for identical inputs, let $\{s_{ij}\}_{i<j}$ be the pairwise STED similarities. Define:
\begin{equation}
    \bar{s} = \frac{2}{n(n-1)} \sum_{i<j} s_{ij}, \qquad \sigma = \text{std}(\{s_{ij}\}_{i<j})
\end{equation}
The \textbf{normalized deviation} is:
\begin{equation}
    \hat{\sigma} = \frac{\sigma}{\sigma_{\max}}, \quad \text{where } \sigma_{\max} = \text{std}(\underbrace{0,\ldots,0}_{\lfloor n(n-1)/4 \rfloor}, \underbrace{1,\ldots,1}_{\lceil n(n-1)/4 \rceil})
\end{equation}
The \textbf{Consistency Score} is:
\begin{equation}
    C(\mathcal{O}) = \left(\frac{1}{1 + 2\hat{\sigma}}\right)^\alpha
\end{equation}
where $\alpha > 0$ is the steepness parameter (default $\alpha = 20$).
\end{definition}

\begin{proposition}[Consistency Score Properties]
\label{prop:consistency-properties}
The consistency score $C(\mathcal{O})$ satisfies:
\begin{enumerate}
    \item[(i)] \textbf{Boundedness}: $C(\mathcal{O}) \in (0, 1]$
    \item[(ii)] \textbf{Perfect consistency}: $C(\mathcal{O}) = 1 \Leftrightarrow \sigma = 0$ (all outputs identical)
    \item[(iii)] \textbf{Monotonicity}: $C$ decreases monotonically with $\hat{\sigma}$
    \item[(iv)] \textbf{Sensitivity control}: Higher $\alpha$ amplifies small deviations
\end{enumerate}
\end{proposition}

\begin{proof}
\textbf{(i)} Since $\hat{\sigma} \geq 0$, we have $1 + 2\hat{\sigma} \geq 1$, thus $(1 + 2\hat{\sigma})^{-\alpha} \leq 1$. Also $(1 + 2\hat{\sigma})^{-\alpha} > 0$ for finite $\hat{\sigma}$.

\textbf{(ii)} $C = 1 \Leftrightarrow (1 + 2\hat{\sigma})^{-\alpha} = 1 \Leftrightarrow \hat{\sigma} = 0 \Leftrightarrow \sigma = 0$.

\textbf{(iii)} Let $f(\hat{\sigma}) = (1 + 2\hat{\sigma})^{-\alpha}$. Then $f'(\hat{\sigma}) = -2\alpha(1 + 2\hat{\sigma})^{-\alpha-1} < 0$ for $\alpha > 0$.

\textbf{(iv)} For small $\hat{\sigma}$, Taylor expansion gives:
\begin{equation}
    C \approx 1 - 2\alpha\hat{\sigma} + O(\hat{\sigma}^2)
\end{equation}
Thus $\frac{\partial C}{\partial \hat{\sigma}}\big|_{\hat{\sigma}=0} = -2\alpha$, showing larger $\alpha$ increases sensitivity.
\end{proof}

\begin{proposition}[Consistency Score Convergence]
\label{prop:convergence}
Let $\{o_i\}_{i=1}^n$ be i.i.d. samples from an LLM's output distribution for a fixed prompt. As $n \to \infty$:
\begin{equation}
    C(\mathcal{O}_n) \xrightarrow{p} C^* = \left(\frac{1}{1 + 2\sigma^*/\sigma_{\max}^*}\right)^\alpha
\end{equation}
where $\sigma^*$ is the population standard deviation of pairwise similarities, and the convergence rate is $O(n^{-\beta})$ with $\beta \geq 0.5$.
\end{proposition}

\begin{proof}
\textbf{Step 1: U-statistic structure.}
The sample variance of pairwise similarities is a U-statistic of order 2:
\begin{equation}
    \hat{\sigma}^2_n = \frac{2}{n(n-1)} \sum_{i<j} (S_{ij} - \bar{s}_n)^2
\end{equation}
By Hoeffding's theorem for U-statistics, $\sqrt{n}(\hat{\sigma}^2_n - \sigma^{*2}) \xrightarrow{d} \mathcal{N}(0, 4\zeta_1^2)$ where $\zeta_1^2 = \text{Var}[\mathbb{E}[(S_{12} - \mu)^2 | o_1]]$.

\textbf{Step 2: Delta method application.}
Since $C = g(\sigma)$ where $g(\sigma) = (1 + 2\sigma/\sigma_{\max})^{-\alpha}$ is continuously differentiable with:
\begin{equation}
    g'(\sigma) = \frac{-2\alpha}{\sigma_{\max}}(1 + 2\sigma/\sigma_{\max})^{-\alpha-1}
\end{equation}
The delta method yields:
\begin{equation}
    \sqrt{n}(C_n - C^*) \xrightarrow{d} \mathcal{N}(0, [g'(\sigma^*)]^2 \cdot 4\zeta_1^2)
\end{equation}

\textbf{Step 3: Rate bound.}
By the CLT for U-statistics: $|C_n - C^*| = O_p(1/\sqrt{n})$.

\textbf{Empirical validation:} Testing across sample sizes $n \in \{5, 10, 20, 50, 100, 200\}$ with 20 trials each, we observe convergence rate $\beta \approx 0.73 > 0.5$ (R$^2$ = 0.93), confirming the theoretical bound is conservative.
\end{proof}

\begin{remark}[Choice of $\alpha = 20$]
The parameter $\alpha = 20$ was selected based on the following considerations:
\begin{enumerate}
    \item \textbf{Empirical observation}: LLM outputs typically exhibit small normalized deviations ($\hat{\sigma} < 0.1$)
    \item \textbf{Discrimination}: With $\alpha = 20$, a 5\% deviation ($\hat{\sigma} = 0.05$) yields $C \approx 0.37$, providing meaningful discrimination
    \item \textbf{Saturation}: For $\hat{\sigma} > 0.2$, the score approaches 0, which is appropriate for highly inconsistent outputs
\end{enumerate}
See Section~\ref{sec:ablation} for ablation studies over $\alpha \in \{5, 10, 15, 20, 25, 30\}$.
\end{remark}

%==============================================================================
\section{Complexity Analysis}
%==============================================================================

\begin{theorem}[STED Computational Complexity]
\label{thm:complexity}
Let $T_1, T_2$ be JSON trees with:
\begin{itemize}
    \item $N = |V_1| + |V_2|$: total node count
    \item $B = \max_{v \in V_1 \cup V_2} |\text{children}(v)|$: maximum branching factor
    \item $D = \max(\text{depth}(T_1), \text{depth}(T_2))$: maximum depth
    \item $S$: maximum string length for embedding computation
\end{itemize}

Then STED computation has:
\begin{equation}
    \text{Time: } O(N \cdot B^3 + N \cdot T_\phi(S)) \qquad \text{Space: } O(B^2 + D + N \cdot d)
\end{equation}
where $T_\phi(S)$ is the embedding computation time and $d$ is embedding dimension.
\end{theorem}

\begin{proof}
\textbf{Time complexity:}
\begin{enumerate}
    \item \textbf{Tree traversal}: Each node visited once: $O(N)$
    \item \textbf{Cost matrix construction}: For internal node with $b$ children, compute $b^2$ pairwise costs. Each cost requires $O(1)$ operations (using cached embeddings). Total: $O(N \cdot B^2)$
    \item \textbf{Hungarian algorithm}: At each internal node, solve $b \times b$ assignment in $O(b^3)$. Summed over all internal nodes: $O(N \cdot B^3)$
    \item \textbf{Embedding computation}: $O(N)$ unique strings, each requiring $O(T_\phi(S))$. With caching, amortized cost: $O(N \cdot T_\phi(S))$
\end{enumerate}

Total: $O(N \cdot B^3 + N \cdot T_\phi(S))$, dominated by Hungarian algorithm for typical JSON with $B \geq 3$.

\textbf{Space complexity:}
\begin{enumerate}
    \item \textbf{Cost matrix}: $O(B^2)$ for largest internal node
    \item \textbf{Recursion stack}: $O(D)$ for tree depth
    \item \textbf{Embedding cache}: $O(N \cdot d)$ for all unique strings
\end{enumerate}

Total: $O(B^2 + D + N \cdot d)$.
\end{proof}

\begin{remark}[Practical Complexity]
For typical production JSON:
\begin{itemize}
    \item $N \in [10, 1000]$, $B \in [2, 20]$, $D \in [2, 10]$
    \item With caching, embedding cost is amortized
    \item Runtime: 10-100ms for moderate JSON, 1-10s for large documents
\end{itemize}

\textbf{Empirical validation:} Memory profiling across $N \in \{10, 25, 50, 100, 200, 500\}$ shows linear scaling with $R^2 > 0.999$. The $N \cdot d$ term (embedding cache with $d = 384$) dominates, contributing $\approx 875$ bytes per node. This confirms the theoretical bound and shows the algorithm is memory-efficient for practical JSON sizes.
\end{remark}

%==============================================================================
\section{Hyperparameter Sensitivity Analysis}
\label{sec:ablation}
%==============================================================================

\begin{proposition}[Structural Threshold Sensitivity]
\label{prop:threshold}
The structural threshold $\theta$ controls the trade-off between structural and content matching:
\begin{itemize}
    \item \textbf{Low $\theta$ (e.g., 0.1)}: Allows content matching between loosely related nodes, increasing flexibility but potentially mismatching unrelated fields
    \item \textbf{High $\theta$ (e.g., 0.7)}: Restricts content matching to structurally similar nodes, preserving semantic alignment but reducing cross-matching opportunities
\end{itemize}
The default $\theta = 0.3$ balances flexibility with semantic coherence.
\end{proposition}

\begin{table}[h]
\centering
\caption{Hyperparameter sensitivity analysis (to be populated with experimental results)}
\label{tab:hyperparams}
\begin{tabular}{lcccc}
\toprule
Parameter & Range & Default & Selection Criterion \\
\midrule
$\alpha$ (consistency steepness) & $[5, 50]$ & 20 & Maximize discrimination in $\hat{\sigma} \in [0, 0.1]$ \\
$\theta$ (structural threshold) & $[0.1, 0.7]$ & 0.3 & Balance flexibility vs. coherence \\
$\lambda$ (path decay) & $[0.5, 1.0]$ & 1.0 & Weight root vs. leaf importance \\
$\lambda_{\text{size}}$ (size penalty) & $[0.05, 0.3]$ & 0.1 & Penalize structural mismatches \\
\bottomrule
\end{tabular}
\end{table}

%==============================================================================
\section{Algorithm Pseudocode}
%==============================================================================

\begin{algorithm}[h]
\caption{STED Similarity Computation}
\label{alg:sted}
\begin{algorithmic}[1]
\REQUIRE Trees $T_1, T_2$, variation type $\in \{\text{struct}, \text{content}, \text{combined}\}$
\ENSURE Similarity score $s \in [0, 1]$
\STATE
\STATE \textbf{function} ComputeSTED$(n_1, n_2)$:
\STATE \hspace{1em} \textbf{if} $(n_1, n_2) \in \text{cache}$ \textbf{then return} $\text{cache}[(n_1, n_2)]$
\STATE
\STATE \hspace{1em} \textbf{if} isLeaf$(n_1) \land$ isLeaf$(n_2)$ \textbf{then}
\STATE \hspace{2em} $\text{cost} \gets \gamma_{\text{type}}(n_1, n_2)$ \COMMENT{Use specified variation type}
\STATE \hspace{1em} \textbf{else if} isLeaf$(n_1) \lor$ isLeaf$(n_2)$ \textbf{then}
\STATE \hspace{2em} $\text{cost} \gets \gamma_{\text{del}}(\text{leaf}) + \sum_{c} \gamma_{\text{ins}}(c)$
\STATE \hspace{1em} \textbf{else}
\STATE \hspace{2em} $C_1 \gets \text{children}(n_1)$, $C_2 \gets \text{children}(n_2)$
\STATE \hspace{2em} $m \gets \max(|C_1|, |C_2|)$
\STATE \hspace{2em} $M \gets \mathbf{0}^{m \times m}$ \COMMENT{Cost matrix}
\STATE \hspace{2em} \textbf{for} $i = 1$ \textbf{to} $m$ \textbf{do}
\STATE \hspace{3em} \textbf{for} $j = 1$ \textbf{to} $m$ \textbf{do}
\STATE \hspace{4em} \textbf{if} $i \leq |C_1| \land j \leq |C_2|$ \textbf{then}
\STATE \hspace{5em} $M[i,j] \gets \text{ComputeSTED}(C_1[i], C_2[j])$
\STATE \hspace{4em} \textbf{else if} $i \leq |C_1|$ \textbf{then}
\STATE \hspace{5em} $M[i,j] \gets \gamma_{\text{del}}(C_1[i])$
\STATE \hspace{4em} \textbf{else if} $j \leq |C_2|$ \textbf{then}
\STATE \hspace{5em} $M[i,j] \gets \gamma_{\text{ins}}(C_2[j])$
\STATE \hspace{2em} $\pi^* \gets \text{Hungarian}(M)$ \COMMENT{Optimal assignment}
\STATE \hspace{2em} $\text{cost} \gets \sum_{(i,j) \in \pi^*} M[i,j]$
\STATE \hspace{2em} $\text{cost} \gets \text{Normalize}(\text{cost}, |C_1|, |C_2|)$
\STATE \hspace{1em} $\text{cache}[(n_1, n_2)] \gets \text{cost}$
\STATE \hspace{1em} \textbf{return} $\text{cost}$
\STATE
\STATE $d \gets \text{ComputeSTED}(\text{root}(T_1), \text{root}(T_2))$
\RETURN $1 - d$
\end{algorithmic}
\end{algorithm}

%==============================================================================
\section{Human Validation Study}
\label{sec:human-validation}
%==============================================================================

To validate that STED captures human judgment of structured output similarity, we conducted a best-match selection study.

\subsection{Study Design}

\textbf{Task.} For each sample, annotators were shown a ground truth JSON output and candidate responses selected by four methods (STED, BERTScore, DeepDiff, TED). Each method selected its ``best match'' to the ground truth based on its similarity score. Annotators identified which candidate was most similar to the ground truth, with method labels anonymized and order randomized.

\textbf{Dataset.} We selected 89 tool-calling samples from the Toucan dataset where the four methods \emph{disagreed} on which response was most similar to ground truth. This focuses evaluation on discriminative cases where method differences matter.

\textbf{Annotators.} Two independent annotators completed the task. Inter-annotator agreement was substantial (Cohen's $\kappa = 0.764$), indicating the task was well-defined.

\subsection{Results}

\begin{table}[h]
\centering
\caption{Best-Match Selection Study Results}
\label{tab:human-validation}
\begin{tabular}{lccc}
\toprule
Method & Win Rate & 95\% CI & $p$-value vs STED \\
\midrule
\textbf{STED} & \textbf{74.7\%} & [64\%, 85\%] & -- \\
BERTScore & 66.9\% & [56\%, 76\%] & 0.31 \\
DeepDiff & 44.4\% & [34\%, 55\%] & $<$0.001*** \\
TED & 24.2\% & [15\%, 34\%] & $<$0.001*** \\
\bottomrule
\end{tabular}
\end{table}

\textbf{Key findings:}
\begin{enumerate}
    \item STED achieves highest alignment with human judgment (74.7\% win rate)
    \item On samples where both annotators agreed ($n=65$), STED reaches \textbf{90.8\%} agreement
    \item STED significantly outperforms DeepDiff ($p < 0.001$) and TED ($p < 0.001$) by McNemar's test
    \item The STED vs BERTScore difference (7.8 percentage points) is consistent across annotators but not statistically significant
\end{enumerate}

\begin{remark}[Interpretation]
The high win rate on agreed items (90.8\%) indicates that when humans are confident about similarity judgments, STED's selections strongly align with human intuition. This validates STED as a reliable proxy for human judgment of structured output similarity, particularly for tool-calling tasks where structural correctness is critical.
\end{remark}

%==============================================================================
\section{Extended Related Work}
%==============================================================================

% Additional baselines to compare against for ICML

\subsection{Learned Edit Distance Methods}

Paaßen et al.~\cite{paassen2018tree} introduced adaptive symbol embeddings for tree edit distance, learning task-specific representations. While effective for classification, their approach requires labeled training data and focuses on discriminative tasks rather than consistency evaluation. STED differs by:
\begin{itemize}
    \item Using pre-trained semantic embeddings (no task-specific training)
    \item Focusing on similarity measurement rather than classification
    \item Providing interpretable cost decomposition (structural vs. content)
\end{itemize}

\subsection{Text Generation Evaluation Metrics}

Recent metrics for text generation evaluation include:
\begin{itemize}
    \item \textbf{BARTScore}~\cite{yuan2021bartscore}: Uses BART likelihood as quality measure; not designed for structured data
    \item \textbf{MoverScore}~\cite{zhao2019moverscore}: Combines contextual embeddings with Earth Mover's Distance; loses structural information when applied to JSON
    \item \textbf{GPTScore}~\cite{fu2023gptscore}: Uses LLM as evaluator; computationally expensive and non-deterministic
\end{itemize}

STED complements these by providing structure-aware evaluation specifically designed for hierarchical outputs.

\subsection{Structured Output Benchmarks}

Recent benchmarks for structured output evaluation:
\begin{itemize}
    \item \textbf{StructEval}~\cite{yang2025structeval}: Benchmarks LLM structural generation capabilities across multiple formats
    \item \textbf{JSONSchemaBench}~\cite{geng2025generating}: Focuses on schema conformance verification
    \item \textbf{BFCL}~\cite{berkeley2024bfcl}: Berkeley Function Calling Leaderboard for tool use evaluation
\end{itemize}

Our consistency framework provides a complementary evaluation dimension focusing on output reliability across repeated generations.

%==============================================================================
% References
%==============================================================================

\begin{thebibliography}{99}

\bibitem{kuhn1955hungarian}
H.~W. Kuhn.
\newblock The Hungarian method for the assignment problem.
\newblock {\em Naval Research Logistics Quarterly}, 2(1-2):83--97, 1955.

\bibitem{munkres1957algorithms}
J.~Munkres.
\newblock Algorithms for the assignment and transportation problems.
\newblock {\em Journal of the SIAM}, 5(1):32--38, 1957.

\bibitem{paassen2018tree}
B.~Paaßen, C.~Gallicchio, A.~Micheli, and B.~Hammer.
\newblock Tree edit distance learning via adaptive symbol embeddings.
\newblock {\em arXiv preprint arXiv:1806.05009}, 2018.

\bibitem{yuan2021bartscore}
W.~Yuan, G.~Neubig, and P.~Liu.
\newblock BARTScore: Evaluating generated text as text generation.
\newblock {\em NeurIPS}, 2021.

\bibitem{zhao2019moverscore}
W.~Zhao, M.~Peyrard, F.~Liu, Y.~Gao, C.~Meyer, and S.~Eger.
\newblock MoverScore: Text generation evaluating with contextualized embeddings and earth mover distance.
\newblock {\em EMNLP}, 2019.

\bibitem{fu2023gptscore}
J.~Fu, S.-K.~Ng, Z.~Jiang, and P.~Liu.
\newblock GPTScore: Evaluate as you desire.
\newblock {\em arXiv preprint arXiv:2302.04166}, 2023.

\bibitem{yang2025structeval}
J.~Yang et al.
\newblock StructEval: Benchmarking LLMs' capabilities to generate structural outputs.
\newblock {\em arXiv preprint arXiv:2505.20139}, 2025.

\bibitem{geng2025generating}
S.~Geng et al.
\newblock Generating structured outputs from language models: Benchmark and studies.
\newblock {\em arXiv preprint arXiv:2501.10868}, 2025.

\bibitem{berkeley2024bfcl}
Berkeley Function Calling Leaderboard.
\newblock \url{https://gorilla.cs.berkeley.edu/leaderboard.html}, 2024.

\end{thebibliography}

\end{document}
